\chapter[CAPÍTULO VIII]{RECOMENDACIONES}

\section{Recomendaciones para la empresa respecto del Proyecto de Mejora / Innovación / Creatividad}

Con base en los resultados y la experiencia obtenida durante el desarrollo e implementación del Sistema Web Help-Desk Centralizado en J\&P Periféricos S.A.C., se formulan las siguientes recomendaciones orientadas a la Gerencia General, la Jefatura de Soporte Técnico y TI, y a los distintos roles operativos del sistema, con el fin de sostener y potenciar los beneficios de la mejora implementada.

\begin{itemize}[leftmargin=1.27cm]
    \item \textbf{Canalización formal de solicitudes:} Se recomienda que las empresas clientes, sucursales y trabajadores utilicen de manera progresiva y sostenida el sistema web como único canal oficial de reporte de incidencias y solicitudes de soporte, reemplazando gradualmente las llamadas telefónicas y los mensajes de WhatsApp sin registro. Para ello, es necesario capacitar a los usuarios en el uso de la plataforma (registro, seguimiento por código PIN, consulta de estados y calificación de la atención), así como establecer una política interna que priorice la gestión de tickets registrados formalmente sobre los canales informales. La consolidación de esta práctica garantizará la trazabilidad del ciclo de vida de cada requerimiento y la disponibilidad de datos para la mejora continua.

    \item \textbf{Mantenimiento y actualización continua del inventario TI multinivel:} Se recomienda mantener actualizado de forma permanente el expediente e inventario de activos TI (hardware, software y licencias) por empresa cliente, sucursal, área y trabajador, registrando cada intervención, componente instalado, actualización o reemplazo. Es fundamental evitar el retorno a planillas dispersas de Excel, asignando la responsabilidad de actualización al personal de soporte y realizando auditorías periódicas del parque informático para asegurar la coherencia entre el inventario registrado y el estado real de los equipos, preservando la trazabilidad histórica que sustenta el diagnóstico y la toma de decisiones técnicas.

    \item \textbf{Gestión ordenada de visitas técnicas \textit{in-situ} y soporte remoto:} Se recomienda respetar las agendas, zonas y sucursales definidas en el módulo de agendamiento, aprovechando el motor de autoasignación por zonas geográficas y el control anti-cruces para agrupar atenciones presenciales de manera eficiente. Asimismo, se debe mantener el flujo de escalamiento estandarizado (derivación por niveles y registro del motivo de reasignación) y coordinar las visitas con información previa del expediente técnico del activo, evitando traslados duplicados, tiempos muertos y reprocesos en la reprogramación de citas.

    \item \textbf{Control periódico de planes, contratos y SLA:} Se recomienda revisar de forma periódica la vigencia de los planes contratados, la cantidad máxima de equipos cubiertos y los tiempos máximos de atención comprometidos, tanto antes como durante la atención de las incidencias. La validación digital de los Acuerdos de Nivel de Servicio (SLA) debe utilizarse de manera sistemática para priorizar las incidencias críticas, evitar la atención de activos fuera de cobertura y garantizar el cumplimiento formal de los compromisos pactados con cada cliente corporativo.

    \item \textbf{Revisión continua de KPIs, respaldos y seguridad:} Se recomienda monitorear de manera continua los indicadores clave de rendimiento (MTTA, MTTR, FCR y CSAT) y las dimensiones de calidad del servicio, utilizando los tableros de control para identificar cuellos de botella y orientar la mejora continua. Asimismo, se deben programar respaldos periódicos de las bases de datos (MySQL, MongoDB y Redis), mantener actualizadas las políticas de seguridad de accesos (JWT y RBAC) conforme a los roles y alcances jerárquicos, y evaluar iterativamente nuevas funcionalidades o ajustes del sistema según las necesidades operativas de J\&P Periféricos S.A.C. y sus empresas clientes.
\end{itemize}

La adopción sostenida de estas recomendaciones permitirá consolidar los beneficios de la mejora implementada, preservar la trazabilidad del servicio de soporte técnico y del inventario de activos TI, y fortalecer la relación de confianza con los clientes corporativos de J\&P Periféricos S.A.C.